\documentclass[11pt]{article}
\usepackage[margin=1in]{geometry}
\usepackage{amsmath,amssymb}
\usepackage{booktabs}
\usepackage{graphicx}
\usepackage[hidelinks]{hyperref}
\usepackage[numbers,sort&compress]{natbib}
\usepackage{xcolor}
\usepackage{lineno}
\linenumbers

\newcommand{\todo}[1]{\textcolor{red}{[TODO: #1]}}
\newcommand{\TD}{TransDecoder}
\newcommand{\PyTD}{PyTransDecoder}

\title{TransDecoder: identification of candidate coding regions within transcript sequences}
\author{Brian J. Haas \and \todo{co-authors}}
\date{}

\begin{document}
\maketitle

\begin{abstract}
\todo{Write after results are final. Points: widely used since 2012 (Trinity, PASA, genome-guided
and long-read transcriptomes); self-training 5th-order Markov model plus GC-calibrated length rule
and optional homology; start-codon refinement; genome projection; new Python implementation
(\PyTD) that reproduces the Perl implementation and runs faster; benchmark on eight eukaryotes
incl.\ a ciliate genetic code and GC extremes against Prodigal, CodAn
\todo{GeneMarkS-T}.}
\end{abstract}

\section{Introduction}
\todo{Context: de novo and genome-guided transcriptome assembly (Trinity \citep{Grabherr2011},
StringTie, long-read isoforms); need to annotate coding regions in transcripts from non-model
organisms; existing approaches: prokaryotic/metagenomic gene finders (Prodigal \citep{Hyatt2010}),
GeneMarkS-T \citep{Tang2015}, CodAn \citep{Nachtigall2021}, ANGEL; homology-based tools.
TransDecoder history: developed with Trinity (\citealp{Haas2013}), used in thousands of studies but
never described in a dedicated publication.}

\section{Methods}
\subsection{Overview}
\TD\ identifies candidate protein-coding regions (CDS) within transcript sequences in two stages
(Fig.~\todo{workflow figure}). \emph{LongOrfs} enumerates open reading frames (ORFs) above a minimum
length. \emph{Predict} trains a coding-sequence model on the longest of these ORFs and uses it,
together with a GC-calibrated length criterion and optional homology evidence, to select the likely
coding ORF(s) in each transcript. It then refines start codons and reports CDS, UTRs and proteins.
An optional third stage projects transcript-level CDS onto a reference genome. Unless stated
otherwise, the algorithm is identical in the Perl implementation (\TD\ v6.1) and in \PyTD\
\todo{version}; Section~\ref{sec:impl} describes their concordance.

\subsection{ORF enumeration}
Each transcript is scanned on the forward strand and, unless the data are declared strand-specific
(\texttt{-S}), on the reverse complement. Stop codons come from the selected genetic code. Twenty-four
named codes are supported, including the ciliate (\emph{Tetrahymena}), \emph{Candida}, Euplotid and
twelve mitochondrial codes. Start codons are ATG in all codes.

In each frame, every stop codon is paired with the most upstream in-frame ATG that follows the
preceding in-frame stop, giving one maximal ORF per stop. ORFs may run off either end of the
transcript, which accommodates incomplete assemblies:
\begin{itemize}
  \item An ORF with no in-frame stop upstream of its first codon is \emph{5$'$ partial}. It begins at
        the first full codon of the transcript.
  \item An ORF that reaches the 3$'$ end without a stop is \emph{3$'$ partial}.
\end{itemize}
ORFs are labelled complete, 5$'$ partial, 3$'$ partial or internal (both). ORFs of at least $m$
codons are retained, where $m=100$ by default; the stop codon, or its absence at a 3$'$ end, counts
as one codon. \texttt{--complete\_orfs\_only} restricts the output to ORFs with both a start and a
stop.

\subsection{Background composition}
Nucleotide frequencies $q(b)$, $b\in\{A,C,G,T\}$, are pooled over all transcripts. Without
\texttt{-S}, counts from both strands are combined, so that $q(A)=q(T)$ and $q(C)=q(G)$. The G+C
content used below is $2\,q(C)$.

\subsection{Training set}
The coding model is trained without supervision on long ORFs, which are very unlikely to arise by
chance:
\begin{enumerate}
  \item Candidates are the $10T$ longest ORFs (by default $T=500$). ORFs with CDS longer than 5{,}000\,nt
        are excluded.
  \item Near-identical proteins are removed greedily, longest first:
        \begin{itemize}
          \item Each candidate protein is cut into consecutive, non-overlapping 5-residue words.
          \item If its distinct words make up less than 30\% of all its words, it is skipped as low
                complexity.
          \item It is discarded if more than 80\% of its distinct words occur in a single
                previously retained protein. Words of retained proteins are indexed at every offset.
        \end{itemize}
  \item The $T$ longest surviving ORFs form the training set. They may be partial and from either
        strand.
\end{enumerate}

\subsection{Coding potential}
Coding sequence is modelled as a three-periodic (inhomogeneous) fifth-order Markov chain
\citep{Borodovsky1993,Parra2000}. Let $N_f(w)$ be the number of training-set occurrences of the
$k$-mer $w$ ($k\le 6$) whose last base lies at codon position $f\in\{0,1,2\}$. The transition
probability of base $x_i$ in frame $f=i \bmod 3$ given the preceding $\min(i,5)$ bases is
\begin{equation}
  P(x_i \mid x_{i-k..i-1}, f) = \frac{N_f(x_{i-k..i}) + 1}{N_{(f-1)\bmod 3}(x_{i-k..i-1}) + 4}.
\end{equation}
Each position contributes the log-odds $\ell_i=\ln\!\left[P(x_i\mid\cdot,f)/q(x_i)\right]$ against the
zeroth-order background. An ORF sequence $s$ is scored as $S(s)=\sum_i \ell_i$, excluding a
terminal stop codon of the selected genetic code.

Six scores are computed for every candidate:
\begin{itemize}
  \item $S_1$ for the ORF in its own frame;
  \item $S_2$, $S_3$ for the same sequence offset by one and two nucleotides;
  \item $S_4$--$S_6$ for the three frames of its reverse complement.
\end{itemize}

\subsection{ORF selection}
An ORF is a candidate coding region if any of the following hold:
\begin{enumerate}
  \item it is supported by homology evidence;
  \item its CDS length is at least $L^\ast(\mathrm{GC})$ (see below);
  \item $S_1>0$ and $S_1>\max(S_2,S_3)$.
\end{enumerate}
Rule~3 compares the reading frame with the two alternative forward frames. An earlier rule also
required $S_1$ to exceed the reverse-complement scores. It was dropped in v5 because coding
sequence on the opposite strand of a gene often scores highly. \todo{cite the 2018 decision-tree
analysis, Supplementary Fig.}

\paragraph{GC-calibrated long-ORF rule.} Long ORFs are kept even when they score poorly. This
recovers genes with atypical composition, such as horizontally transferred genes or genes in
transcriptomes containing several organisms. $L^\ast$ is the 0.999 quantile of the lengths of ORFs
that LongOrfs extracts from random sequence of the same G+C content. The random sequences are
i.i.d., scanned on both strands, with partial ORFs allowed and a 300-nt minimum. $L^\ast$ was
estimated by simulation in 5\% GC steps between 25\% and 80\%, with at least 85{,}000 random ORFs per
level (Supplementary Table~\todo{S1}). It rises from 519\,nt at 25\% GC to 2{,}436\,nt at 80\%. The
rule is disabled above 80\% GC. A fixed length can be set instead
(\texttt{--retain\_long\_orfs\_mode strict}).

\paragraph{Homology.} Optionally, the proteins of all candidate ORFs can be searched against a
protein database with DIAMOND \citep{Buchfink2021} or BLASTP, and against Pfam \citep{Mistry2021}
with HMMER \citep{Eddy2011}. ORFs with hits are retained and ranked first.

\paragraph{Ranking and overlap resolution.} Within each transcript, candidates are ranked by:
\begin{enumerate}
  \item number of homology hits;
  \item $S_1$;
  \item length.
\end{enumerate}
They are then accepted greedily. A candidate is rejected if it overlaps an accepted ORF on either
strand by more than 10\% of either ORF's length; homology-supported ORFs are always accepted.
Several non-overlapping ORFs per transcript can therefore be reported, as in polycistronic or
chimeric assemblies. \texttt{--single\_best\_only} keeps a single ORF per transcript.

\subsection{Start-codon refinement}
\label{sec:start}
A 5$'$-partial ORF begins at the first in-frame codon of the transcript rather than at an ATG. If
the transcript is complete, the true start lies downstream. \TD\ learns the sequence context of
translation start sites from the training set:
\begin{itemize}
  \item Each feature spans 20\,nt upstream of an ATG, the ATG, and 10\,nt downstream.
  \item Positive features are the start codons of training ORFs.
  \item Negative features are all ATGs downstream of those starts.
\end{itemize}
Position weight matrices (pseudocount 0.1) are estimated for positives and negatives. A candidate is
scored as the log-likelihood ratio over a sub-window of $u$ upstream and $d$ downstream positions,
divided by $u+d$.

The procedure has three steps:
\begin{enumerate}
  \item Noisy positives are depleted iteratively, starting from a random 30\% subset.
  \item The sub-window $(u,d)$, with $1\le u\le20$ and $1\le d\le10$, is chosen by maximising the
        cross-validated area under the ROC curve. The decision threshold is chosen by maximising
        the F1 score.
  \item For each ORF that does not begin with ATG, the in-frame ATGs within
        $\max(30\,\mathrm{nt}, 15\%\ \text{of ORF length})$ of its 5$'$ end are scored. The ORF is
        moved to the highest-scoring ATG at or above the threshold, and relabelled complete (or
        3$'$ partial).
\end{enumerate}
\todo{Update with the stabilised procedure if adopted (analyses/start\_refinement\_stability). Note the
2026 fix: earlier releases learned the threshold on length-normalised scores but applied it to
unnormalised scores on a window shifted by one nucleotide.}

\subsection{Outputs and genome projection}
\TD\ writes GFF3 (gene, mRNA, CDS, UTR), BED, and CDS and peptide FASTA files. The record headers
carry ORF type, score and homology hits.

For transcripts reconstructed from a genome alignment (StringTie, Cufflinks, PASA, GMAP, or
long-read isoform pipelines), \texttt{cdna\_alignment\_orf\_to\_genome\_orf.pl} projects the
transcript CDS onto the genome through the exon structure, giving genome-coordinate gene models in
GFF3. ORFs from different isoforms of a gene keep a common gene identifier. ORFs predicted on the
opposite strand of a multi-exon transcript are discarded. For a single-exon transcript, such an ORF
reverses the transcript's orientation.

\subsection{Implementations}
\label{sec:impl}
\TD\ is implemented in Perl (\url{https://github.com/TransDecoder/TransDecoder}). Version~6 adds a
single wrapper that runs genome-based cDNA extraction, LongOrfs, optional DIAMOND/BLASTP and
HMMER searches, Predict, and genome projection.

\PyTD\ is a pure-Python reimplementation with Biopython as its only dependency. It reproduces the
Perl algorithm, including the ordering of ties. On five test data sets the two implementations
produce identical ORFs, coding scores and selections; they differ only in start codons chosen by the
randomised PWM training (Section~\ref{sec:start}; Supplementary Table~\todo{S2}). \PyTD\ runs
\todo{2.7}$\times$ faster \todo{(benchmark runtime table)}.

\subsection{Benchmark data}
Reference coding regions were taken from the RefSeq \citep{OLeary2016} transcript records (\texttt{rna.gbff}) of eight
eukaryotes (Table~\ref{tab:orgs}):
\begin{itemize}
  \item the four organisms of the original 2018 evaluation;
  \item human, restricted to MANE Select transcripts \citep{Morales2022};
  \item \emph{Tetrahymena thermophila}, which uses the ciliate genetic code: TAA and TAG encode
        glutamine, and TGA is the only stop;
  \item two compositional extremes: \emph{Plasmodium falciparum} (\todo{x}\% GC) and
        \emph{Chlamydomonas reinhardtii} (\todo{y}\% GC).
\end{itemize}
Curated (NM\_) records were used where available; model (XM\_) records were used for the last three.

A transcript was eligible if its single annotated CDS:
\begin{itemize}
  \item was complete and contiguous;
  \item began with ATG and ended with a stop of the record's translation table;
  \item had no internal stop, frameshift, translational exception or ribosomal slippage;
  \item translated to the annotated protein.
\end{itemize}
For each gene the eligible transcript with the longest CDS was kept. UTRs were trimmed to at most
250\,nt, as in the 2018 evaluation.

Each truth set was used to derive four inputs (all generated with fixed seeds):
\begin{enumerate}
  \item \emph{complete}: the transcripts as given;
  \item \emph{fragments}: each transcript is cut once within its CDS at a fraction drawn from
        $\mathcal N(0.5,0.1)$. Cuts 3$'$ of the CDS midpoint keep the 3$'$ piece, which retains the
        stop but not the start; other cuts keep the 5$'$ piece, which retains the start but not the
        stop;
  \item \emph{unstranded} versions of (1) and (2), in which a random half of the sequences are
        reverse-complemented;
  \item each of the above with one random negative sequence per transcript added. A negative has
        the transcript's length and G+C content.
\end{enumerate}

\subsection{Compared methods}
The following were run with default parameters unless stated:
\begin{itemize}
  \item \PyTD, with and without start refinement;
  \item Perl \TD\ (current release, v6.0.0 and v5.7.1);
  \item Prodigal 2.6.3 \citep{Hyatt2010}, in single mode with a full motif scan (\texttt{-n}) and
        the organism's translation table;
  \item CodAn 1.1 \citep{Nachtigall2021}, with the pre-trained model of the closest clade (full or
        partial model to match the input);
  \item \todo{GeneMarkS-T \citep{Tang2015}};
  \item two baselines: the longest ORF per transcript, and all ORFs of at least 100 codons.
\end{itemize}
Inputs labelled strand-specific were analysed on the forward strand only (\TD\ \texttt{-S};
Prodigal forward-strand calls; CodAn \texttt{--strand plus}). Unstranded inputs were analysed on both
strands.

\subsection{Accuracy measures}
Following the 2018 evaluation, a prediction matched a reference CDS on the same transcript and strand
when its 3$'$ end was within 2\,nt of the reference 3$'$ end (stop-level match). It was an exact match
when its 5$'$ end was also within 2\,nt. Each reference CDS could be matched at most once.

We report:
\begin{itemize}
  \item sensitivity (matched references / references);
  \item precision (matched predictions / all predictions, including predictions on random
        negatives);
  \item their harmonic mean (F1).
\end{itemize}
\todo{Runtime measured on \ldots}

\subsection{Ribosome profiling support}
GM12878 ribosome footprints \citep{Cenik2015} (GEO GSE65912; SRR1803147, SRR1803148) were processed
as follows:
\begin{itemize}
  \item adapter-trimmed with cutadapt, keeping 25--34~nt;
  \item depleted of rRNA, mitochondrial rRNA and tRNA reads with bowtie2;
  \item aligned end-to-end with STAR to GRCh38 with GENCODE v50 splice junctions, keeping unique
  alignments only (11.5 million footprints).
\end{itemize}
P-site offsets were calibrated per read length on MANE Select CDSs. For each length we chose the
5$'$ offset that maximised the frame-0 fraction, and used only lengths with a frame-0 fraction
$\ge$0.5 (27, 29 and 30~nt). For each predicted CDS, P-sites were counted by frame in spliced CDS
coordinates, excluding the first and last codons. A CDS was called supported when it had at least
ten P-sites, frame~0 was the most frequent frame, and the binomial probability of the frame-0 count
under $p=1/3$ was below 0.01.

\subsection{Mass-spectrometry support}
We used the published MaxQuant peptide table for GM12878 \citep{Sinitcyn2023} (MassIVE
MSV000086944), with peptides already filtered at 1\% PSM and protein FDR. We removed reverse and
contaminant entries, then matched peptides exactly (I and L treated as equivalent, Aho--Corasick)
to the predicted proteins, GENCODE v50 translations and the UniProt human proteome. A peptide was
called non-reference when it occurred in neither reference set.


\section{Results}
\todo{(1) Perl/Python concordance and runtime; (2) benchmark accuracy across organisms and settings;
(3) effect of non-standard genetic code; (4) GC extremes and the long-ORF rule; (5) start-codon
refinement; (6) homology integration; (7) application examples (Trinity assembly + BUSCO;
long-read isoforms + genome projection).}
\subsection{Translation and protein evidence for predicted ORFs in GM12878}
% Draft (2026-10-21). Numbers: genome_annotation/results/gm12878_riboseq_*.tsv, gm12878_ms_*.tsv;
% notes/06_gm12878_three_assemblies.md.

We tested the GM12878 predictions against two independent experimental data sets from the same cell
line: ribosome profiling \citep{Cenik2015} and a deep, six-protease proteome \citep{Sinitcyn2023}.
Each distinct predicted CDS was classed by its relation to GENCODE: exact CDS match, compatible
(a correct partial CDS), other overlap with an annotated CDS, or unannotated (no overlap with any
annotated CDS).

\paragraph{Ribosome footprints.}
Predictions matching annotated CDSs were broadly covered by footprints. Of exact-match predictions,
83--89\% had at least ten P-sites, and 53--59\% passed a test for frame-0 enrichment. That test is
conservative here, because the RNase~A/T1 digestion used for these libraries gives weak 3-nt
periodicity (frame-0 fraction about 0.55 for the read lengths used). Unannotated predictions were
almost never covered: 4.3\%, 1.0\% and 0.6\% (LRAA, StringTie, Trinity), with 0.2--1.5\% passing
the frame test. The difference held when we compared transcripts with similar RNA abundance. At
5--20 TPM, 89\% of exact-match predictions were covered against 3.4\% of unannotated ones, so low
expression does not explain the lack of footprints.

\paragraph{Mass-spectrometry peptides.}
Of 758,603 peptides identified at 1\% FDR, 92--94\% of exact-match predictions of at least 100~aa
carried at least one peptide, nearly all of them two or more. Short exact-match predictions
(\textless100~aa) had 81--87\% coverage. Only 10--25\% of unannotated predictions shared any
peptide with the identified set, and these were almost entirely paralogous or pseudogene copies of
known protein-coding genes (for example \textit{NCF1B/C}, \textit{HERC2P3} and processed
ribosomal-protein pseudogenes). There, peptides cannot distinguish the copy from its parent gene.
Peptides absent from both GENCODE and UniProt that matched unannotated predictions were nearly all
GM12878 variant-allele peptides of parent genes mapping to duplicated copies, for example a
\textit{PWP2} copy on 21p and the \textit{NAPSB} pseudogene. The one credible exception was
\textit{C3orf86P}, supported by three independent peptides from three proteases. It was annotated
as protein-coding in the Ensembl release used for the original search and is a transcribed unitary
pseudogene in GENCODE v50. Because the original search database did not include our predictions,
this analysis confirms known proteins but cannot discover new ones.

\paragraph{Interpretation.}
Ribosome profiling and proteomics agree. Predictions at annotated coding loci are translated and
detected as protein, while the short unannotated ORFs gained at a 50-aa minimum show no evidence of
translation or protein in this cell line. Their apparent protein support comes from paralogs and
pseudogenes. Such predictions are best treated as candidates for evidence-based confirmation, not
as proteins.
\todo{Results for the homology-restriction series (clade-excluded Swiss-Prot); link to the
minimum-length sweep; candidate neo-ORFs supported only by transposon homology.}


\section{Discussion}
\todo{}

\section*{Availability}
\TD\ (Perl): \url{https://github.com/TransDecoder/TransDecoder}.
\PyTD: \url{https://github.com/TransDecoder/PyTransDecoder}.
Benchmark data, scripts and analyses: \url{https://github.com/TransDecoder/TransDecoderPaper}.

\bibliographystyle{unsrtnat}
\bibliography{references}
\end{document}
